\chapter{Conclusiones}

En este trabajo se analizó el uso de la señal ISDB-Tb como iluminador de oportunidad en un sistema de radar pasivo biestático. Con ese objetivo, se estudió la estructura del estándar y, a partir de ello, se desarrolló una cadena de procesamiento.

A partir de las simulaciones realizadas, se observó que con la reconstrucción de la referencia se obtiene una mejora en el desempeño de la cadena, ya que se obtuvo una reducción entre 3 y 4 dB en el piso de ruido en el bloque de detección, cabe aclarar que las aproximaciones dependen directamente de las configuraciones del bloque CFAR y del bloque de ventaneo por lo que pueden variar para un mismo escenario y señal. En este sentido, esto se refleja en una mayor sensibilidad y en mejores condiciones para la detección de blancos.

Sin embargo, también se notó que con la mejora de ruido también se hicieron más evidentes ciertas componentes determinísticas dadas por el carácter cíclico y periódico de la señal. Esto no necesariamente invalida la utilización de este enfoque, simplemente hace notoria la necesidad de un postprocesamiento para limitar las detecciones espurias.

Con esto, se considera que se cumplieron los objetivos planteados al inicio del trabajo, ya que se logró analizar el estándar ISDB-Tb con el nivel de detalle correcto para el contexto y, por otro, se desarrolló una cadena de procesamiento completa que permitió analizar, en un entorno controlado, el efecto de reconstruir la señal de referencia sobre el desempeño del sistema.

\section{Trabajo futuro}

Las posibles mejoras y adiciones para continuar el trabajo presentado es que excedieron el marco del mismo, son la validación de la cadena utilizando señales reales para poder incorporar de manera más realista los efectos del canal, errores de sincronización y otras perturbaciones que no fueron contempladas en el modelo utilizado. También sería conveniente seguir con el análisis de la mitigación de las componentes determinísticas con estrategias como la del filtro recíproco.

